%!TEX root = TheoreticalPhysicsC.tex

\section{Lecture: Spin}
The angular momentum operators satisfy the following commutation relations:
\begin{equation}\label{eq:amcr}
	[\widehat{L}_1,\widehat{L}_2] = i\hbar \widehat{L}_3, \quad [\widehat{L}_2,\widehat{L}_3] = i\hbar \widehat{L}_1\quad\text{and}\quad [\widehat{L}_3,\widehat{L}_1] = i\hbar \widehat{L}_2.
\end{equation}
It follows that
\begin{equation}
	[\widehat{\mathbf{L}}^2, \widehat{L}_j]=0, \quad \text{for all $j=1,2,3$,}
\end{equation}
where
\begin{equation}
	\widehat{\mathbf{L}}^2 \equiv \widehat{L}_1^2+\widehat{L}_2^2+\widehat{L}_3^2.
\end{equation}
We have seen that the orbital angular momentum operators $\widehat{\mathbf{L}} \equiv \widehat{\mathbf{r}}\times \widehat{\mathbf{p}}$ satisfy these commutators. One says that $\widehat{\mathbf{r}}\times \widehat{\mathbf{p}}$ supplies us with a \emph{representation} of these commutators. An important property of this representation is that the total angular momentum is \emph{quantised}, in the sense that the eigenvalues of the total angular momentum operator $\widehat{\mathbf{L}}^2$ can take only the values $\hbar^2 \ell(\ell+1)$, where $\ell = 0, \frac12, 1, \ldots$. For the orbital angular momentum operators $\widehat{\mathbf{r}}\times \widehat{\mathbf{p}}$, however, only integer values show up: $\ell = 0,1,2,\ldots$.

Every representation comes with a Hilbert space on which the operators act; the only one we have discussed in detail so far is $\mathcal{H}=L^2(\mathbb{R}^3)$.

A fascinating question remains open: are there quantum systems which exhibit half-integer eigenvalues? To answer this question in the affirmative we need a Hilbert space $\mathcal{H}$ and three operators $\widehat{S}_j:\mathcal{H}\rightarrow \mathcal{H}$, $j=1,2,3$, which satisfy the relations (\ref{eq:amcr}):
\begin{equation}\label{eq:amcr-spin}
	\eqbox{[\widehat{S}_1,\widehat{S}_2] = i\hbar \widehat{S}_3, \quad [\widehat{S}_2,\widehat{S}_3] = i\hbar \widehat{S}_1\quad\text{and}\quad [\widehat{S}_3,\widehat{S}_1] = i\hbar \widehat{S}_2.}
\end{equation}
If such a Hilbert space and such operators exist, then we can at least establish that there is no \emph{mathematical} obstruction preventing the existence of quantum systems with half-integer total angular momentum. Even if such a quantum system can exist mathematically, it is still a further step to establish whether such a system can occur in nature.

It turns out that quantum systems with half-integer eigenvalues for the total angular momentum operator do indeed exist. We now investigate the case $\ell = \frac12$. The associated Hilbert space is given by $\mathcal{H} \equiv \mathbb{C}^2$. That we now call the operators $\widehat{S}_j$ instead of $\widehat{L}_j$ is nothing more than a renaming. To this end we define the matrices
\begin{equation}
	\widehat{S}_1 = \frac{\hbar}{2} \begin{pmatrix}
		0 & 1 \\ 1 & 0
	\end{pmatrix}, \quad \widehat{S}_2 = \frac{\hbar}{2} \begin{pmatrix}
		0 & -i \\ i & 0
	\end{pmatrix} \quad \text{and} \quad \widehat{S}_3 = \frac{\hbar}{2} \begin{pmatrix}
		1 & 0 \\ 0 & -1
	\end{pmatrix}.
\end{equation}
Up to the factor $\frac{\hbar}{2}$ these are the \emph{Pauli matrices} $\sigma_1,\sigma_2,\sigma_3$. It remains to check whether they satisfy (\ref{eq:amcr-spin}):
\begin{equation}
	[\widehat{S}_1,\widehat{S}_2] = \frac{\hbar^2}{4}\left(\sigma_1\sigma_2-\sigma_2\sigma_1\right) = \frac{\hbar^2}{4}\cdot 2i \begin{pmatrix}
		1 & 0 \\ 0 & -1
	\end{pmatrix} = i\hbar\, \widehat{S}_3.
\end{equation}
The remaining two relations follow the same pattern.

\begin{historical}
	This answers only the \emph{mathematical} question: one can define wild Hilbert spaces which satisfy every condition and nevertheless describe no system in nature. With spin we are in luck -- in the Stern-Gerlach experiment, which split a beam of silver atoms into exactly two partial beams, such a system was found: the electron possesses an \emph{internal} degree of freedom with total angular momentum $\frac12$.
\end{historical}

The total angular momentum operator $\widehat{\mathbf{S}}^2$ reads
\begin{equation}
	\widehat{\mathbf{S}}^2 = \frac{3\hbar^2}{4} \mathbb{I}.
\end{equation}
The calculation fits on a single line, since $\sigma_1^2 = \sigma_2^2 = \sigma_3^2 = \mathbb{I}$. This operator possesses two (degenerate) eigenvalues, namely $\frac{3\hbar^2}{4}$ and $\frac{3\hbar^2}{4}$. Comparing the action of $\widehat{\mathbf{S}}^2$ with
\begin{equation}
	\widehat{\mathbf{S}}^2|\ell,m\rangle = \hbar^2\ell(\ell+1)|\ell,m\rangle
\end{equation}
we establish that $\ell=\frac12$ must hold, since $\frac34 = \frac12\left(\frac12+1\right)$.

That $[\widehat{\mathbf{S}}^2,\widehat{S}_j]=0$ holds is trivial here: a multiple of the identity operator commutes with everything. With respect to this representation $\widehat{S}_3$ is already diagonal. That is to say, the eigenbasis of $\widehat{S}_3$ may be identified with $|\ell,m\rangle$:
\begin{equation}
	|\ell=\tfrac12,m=\tfrac12\rangle \equiv \begin{pmatrix}
		1 \\ 0
	\end{pmatrix}\quad\text{and}\quad 
	|\ell=\tfrac12,m=-\tfrac12\rangle \equiv\begin{pmatrix}
		0 \\ 1
	\end{pmatrix}.
\end{equation}
These are also written, for short, as $|\uparrow\rangle$ and $|\downarrow\rangle$.

Representations for $\ell=\frac32,\frac52,\ldots$ can be written down just as well, but nobody has yet discovered \emph{fundamental} particles with $\ell=\frac32$. Such angular momenta can, however, be \emph{built up}: a system with $\ell=\frac12$ brought together with one with $\ell=1$ looks, from afar, like $\ell=\frac32$. With that we have arrived at composite systems.

\section*{Composite Systems}
Often, in modelling a quantum mechanical system, we have to build the system up out of several subsystems: an electron and a proton then become a hydrogen atom, and so on; or we refine our description by bringing in further degrees of freedom: an ``electron without spin'', together with a description of an isolated spin degree of freedom, becomes a theory of electrons that includes the spin.

In all these cases we proceed by the rule that every kind of quantum mechanical
system is assigned a Hilbert space. So let $\mathcal{H}_1$ be the Hilbert space of the kind of system ``1'' and likewise $\mathcal{H}_2$. The question is then: how, and in which Hilbert space ``$\mathcal{H}_{\text{$1$\&$2$}}$'', should we describe the kind of system ``1 \& 2'', made up of a pair of particles/degrees of freedom of kinds 1 and 2? The answer is given with the help of the \emph{tensor product} $\mathcal{H}_1\otimes \mathcal{H}_2$ of $\mathcal{H}_1$ and $\mathcal{H}_2$.

Out of two Hilbert spaces a third can be built in several ways; only one of them is the physically correct one. The requirement is this: the subsystems must be preparable \emph{independently} of one another -- it must make no difference whether a second hydrogen atom is standing next to the first. Every combination of states must therefore correspond to a vector in $\mathcal{H}_{\text{$1$\&$2$}}$, and from that it follows that
\begin{equation}
	\dim\left(\mathcal{H}_{\text{$1$\&$2$}}\right) = \dim(\mathcal{H}_1)\cdot \dim(\mathcal{H}_2):
\end{equation}
The dimensions \emph{multiply}, they do not add -- as with classical bits, where three bits possess eight states.

\begin{center}
\parbox[c][5cm][c]{14cm}{\begin{tcolorbox}[colback=banana,arc=0cm,frame hidden,enhanced]{The tensor product $\mathcal{H}_1\otimes \mathcal{H}_2$ of two Hilbert spaces $\mathcal{H}_1$ and $\mathcal{H}_2$ is \emph{the} Hilbert space which contains the vectors $\phi_1\otimes \phi_2$ with $\phi_1\in \mathcal{H}_1$ and $\phi_2\in \mathcal{H}_2$, and is generated by them, and in which the computational rule 
\begin{equation}
	\langle \phi_1\otimes \phi_2|\psi_1\otimes\psi_2\rangle \equiv \langle\phi_1| \psi_1\rangle\langle\phi_2| \psi_2\rangle
\end{equation}
holds for inner products.
}\end{tcolorbox}}
\end{center}

Let us next consider the tensor product of operators. For $A_1:\mathcal{H}_1 \rightarrow \mathcal{H}_1$ and $A_2:\mathcal{H}_2 \rightarrow \mathcal{H}_2$ we define the operator $A_1\otimes A_2:\mathcal{H}_1 \otimes\mathcal{H}_2 \rightarrow \mathcal{H}_1 \otimes\mathcal{H}_2$ by the equation
\begin{equation}
	\eqbox{(A_1\otimes A_2)(\phi_1\otimes \phi_2) \equiv (A_1\phi_1)\otimes (A_2\phi_2).}
\end{equation}
Since $A_1\otimes A_2$ is meant to be a linear operator, it is thereby defined on all
linear combinations of product vectors as well.

A tensor product of operators can of course also be given by its matrix elements. Directly from the definition it follows that
\begin{equation}
	\begin{split}
			[A_1\otimes A_2]_{jj',kk'} &\equiv \langle \phi_j\otimes \phi_{j'}|\left(A_1\otimes A_2\right) \phi_k\otimes\phi_{k'}\rangle \\
			&= \langle \phi_j\otimes \phi_{j'}| (A_1\phi_k)\otimes (A_2\phi_{k'})\rangle \\
			&= \langle \phi_j|A_1\phi_k\rangle \langle \phi_{j'}|A_2\phi_{k'}\rangle \\
			&= [A_1]_{j,k}[A_2]_{j',k'}.
	\end{split}
\end{equation}
The indices on the product matrix are thus pairs of indices. The tensor product of a $3\times 3$ matrix with a $7\times 7$ matrix therefore yields a $21\times 21$ matrix. For the numerical handling of such products it is necessary to turn the pairs of indices back into simple indices, that is, to number all the possible pairs consecutively in some way.

Computational rules for the tensor product of operators usually follow most easily from the definition. Alongside the linearity in each factor, which holds for all tensor products (those of vectors included), we obtain:
\begin{align}
	(A_1\otimes A_2)(B_1\otimes B_2) &= A_1B_1\otimes A_2B_2 \\
	(A_1\otimes A_2)^\dag &= A_1^\dag \otimes A_2^\dag \\
	\text{$A_1\ge 0$ and $A_2\ge 0$} &\Rightarrow A_1\otimes A_2 \ge 0 \\
	\|A_1\otimes A_2\| &= \|A_1\|\| A_2\| \\
	\tr(A_1\otimes A_2) &= \tr(A_1)\tr(A_2).
\end{align}

Finally it should be mentioned that everything said here about the tensor product of two factors carries over at once to finitely many factors.

\section*{Dynamics of Composite Systems}
If both subsystems run independently of one another, then the time evolution of the total system is given by $U_t = U^{(1)}_t\otimes U^{(2)}_t$. With the Hamiltonian, by contrast, there is a trap lurking: one would guess it to be the product $\widehat{H}_1\otimes\widehat{H}_2$ -- and that is wrong. The correct answer is
\begin{equation}\label{eq:H-tensor}
	\eqbox{\widehat{H} = \widehat{H}_1\otimes \mathbb{I} + \mathbb{I}\otimes \widehat{H}_2.}
\end{equation}
The reason is the product rule: differentiating $U_t$ with respect to time gives
\begin{equation}
	\frac{d}{dt}\left(U^{(1)}_t\otimes U^{(2)}_t\right) = \left(\frac{dU^{(1)}_t}{dt}\right)\otimes U^{(2)}_t + U^{(1)}_t\otimes\left(\frac{dU^{(2)}_t}{dt}\right),
\end{equation}
and from this, with $i\hbar\frac{dU^{(j)}_t}{dt} = \widehat{H}_j U^{(j)}_t$ and $i\hbar\frac{dU_t}{dt} = \widehat{H}U_t$, one reads off precisely (\ref{eq:H-tensor}). Energies add, they do not multiply. Once the partial problems $\widehat{H}_1|E_j\rangle = E_j|E_j\rangle$ and $\widehat{H}_2|E_k\rangle = E_k|E_k\rangle$ are solved, $\widehat{H}$ is diagonalised too:
\begin{equation}
	\begin{split}
		\widehat{H}\left(|E_j\rangle\otimes |E_k\rangle\right) &= \left(\widehat{H}_1|E_j\rangle\right)\otimes |E_k\rangle + |E_j\rangle \otimes \left(\widehat{H}_2|E_k\rangle\right) \\
		&= (E_j+E_k)\, |E_j\rangle\otimes |E_k\rangle.
	\end{split}
\end{equation}
\begin{exercise}
	Show that these product vectors form a \emph{complete} eigenbasis of $\widehat{H}$.
\end{exercise}

\section*{Entanglement}
One last remark, and it is the most consequential one. For $\mathcal{H} = \mathbb{C}^2\otimes\mathbb{C}^2$ with orthonormal bases $\{\psi_1,\psi_2\}$, $\{\phi_1,\phi_2\}$ every pure state is a linear combination of the four product vectors $\psi_j\otimes\phi_k$ -- but not every such combination is itself a product. An example is the Einstein-Podolsky-Rosen state ($\psi_1=\phi_1=|\uparrow\rangle$, $\psi_2=\phi_2=|\downarrow\rangle$):
\begin{equation}
	|\text{EPR}\rangle = \frac{1}{\sqrt{2}}\left(\psi_1\otimes\phi_1 + \psi_2\otimes\phi_2\right) = \frac{1}{\sqrt{2}}\left(|\uparrow\uparrow\rangle + |\downarrow\downarrow\rangle\right).
\end{equation}
It cannot be written as $\psi\otimes\phi$: such a product $(a_1\psi_1+a_2\psi_2)\otimes(b_1\phi_1+b_2\phi_2)$ has the coefficients $a_jb_k$, and from $a_1b_2 = a_2b_1 = 0$ it follows that $a_1b_1 = 0$ or $a_2b_2 = 0$ -- both cannot be $\tfrac{1}{\sqrt{2}}$. It therefore cannot be read as an independent preparation of two systems -- although it was precisely independent preparability that led us to the tensor product. The tensor product thus delivers more than we asked of it.

% fig_vl17_bloch.tex — Spin 1/2: die Blochkugel eines einzelnen Spins, und
% was aus dem Blochpfeil wird, sobald der Spin zu einem verschränkten Paar
% gehört.
%
% Follows the conventions of fig_vl11_kasten (the reference plate).
%
%   * geometry comes from the formulae of the lecture, never freehand.
%
%     Panel (a).  Every pure state of H = C^2 is, up to a global phase,
%         |psi(theta,phi)> = cos(theta/2)|up> + e^{i phi} sin(theta/2)|down>,
%     and with the S_j of the lecture its vector of expectation values is
%         (2/hbar)(<S_1>,<S_2>,<S_3>)
%             = (sin th cos ph, sin th sin ph, cos th),
%     a unit vector.  The panel therefore draws the unit sphere of that
%     vector, and every marked point sits where the formula puts it:
%     |up>,|down> at th = 0,pi — antipodal, not perpendicular, because the
%     half angle th/2 turns the 90 degrees of orthogonality in C^2 into
%     180 degrees in R^3; |+>,|-> at th = pi/2, ph = 0,pi;  |+i>,|-i> at
%     th = pi/2, ph = +-pi/2.  The projection is a genuine orthographic
%     one, camera azimuth psi = 25 deg and elevation eps with sin eps = 0.3:
%         x^ -> (-0.42262, -0.27189)
%         y^ -> ( 0.90631, -0.12679)
%         z^ -> ( 0.00000,  0.95394)
%     Consequences that are therefore forced rather than drawn by eye: the
%     equator projects to an ellipse of axis ratio 0.3; its *near* half is
%     the lower one, hence solid, the far half dashed; and the poles land
%     at 0.954 R, just inside the outline, which is what an orthographic
%     projection of a sphere actually does.  \tpcbx/\tpcby implement the
%     projection, so no coordinate in the panel was typed by hand.
%
%     Panel (b).  The one-parameter family
%         |psi_a> = cos(a)|up up> + sin(a)|down down>
%     runs from a product state through the EPR state of the lecture
%     (a = pi/4) to another product state.  Both curves are exact:
%         (2/hbar)|<S (x) 1>| = |cos^2 a - sin^2 a| = |cos 2a|,
%             the length of spin 1's arrow in panel (a) — for this family
%             <S_1 (x) 1> = <S_2 (x) 1> = 0 and <S_3 (x) 1> = (hbar/2)cos 2a,
%             all three computable with the tensor-product rules of the
%             lecture alone;
%         2|c_11 c_22 - c_12 c_21| = 2|cos a sin a| = |sin 2a|,
%             exactly the obstruction the lecture computes: a product needs
%             a_1 b_2 = a_2 b_1 = 0, i.e. a vanishing determinant of the
%             coefficient matrix.
%     Since cos^2 + sin^2 = 1 the two readings trade off exactly, and the
%     crossing points at a = pi/8, 3pi/8 are at 1/sqrt(2).  At the EPR
%     point the arrow of panel (a) has length zero: it has retreated from
%     the surface of the ball to its centre.  That is the whole content of
%     the plate — the two panels are one picture, read twice.
%
%   * colour carries the observable and nothing else: eigenstates of S_3
%     navy, of S_1 blue, of S_2 teal (the ordered ramp of `tpclevels`);
%     the one state that is not a basis state, the general |psi>, is
%     orange, and so is the entanglement curve in (b) — the blue curve in
%     (b) is the length of a blue-and-navy object, the arrow.  Axes,
%     equator and the EPR rule are neutral grey: reference structure.
%   * near/far is coded by fill, not by colour: a point of the near
%     hemisphere is a filled disc, one of the far hemisphere an open
%     circle, and the far halves of the axes are dashed.  This survives
%     greyscale printing and makes the antipodal pairs read as pairs.
%   * (b) separates its curves by dash pattern as well as by colour, for
%     the same reason.
%
% Panel (a) the state space of one spin 1/2.
% Panel (b) what becomes of that arrow when the spin is one of a pair.

% Orthographic projection of the sphere of radius #1, see the header.
% #2 = polar angle theta, #3 = azimuth phi, both in degrees.
\def\tpcbx#1#2#3{#1*(-0.42262*sin(#2)*cos(#3)+0.90631*sin(#2)*sin(#3))}
\def\tpcby#1#2#3{#1*(-0.27189*sin(#2)*cos(#3)-0.12679*sin(#2)*sin(#3)%
                     +0.95394*cos(#2))}
% \tpcbc{name}{radius}{theta}{phi} — name a point on that sphere.
\newcommand{\tpcbc}[4]{%
  \coordinate (#1) at ({\tpcbx{#2}{#3}{#4}},{\tpcby{#2}{#3}{#4}});}
\def\tpcRb{1.7}

\begin{center}
\begin{tikzpicture}[
    % near hemisphere: filled disc;  far hemisphere: open circle
    sdot/.style ={circle, inner sep=0pt, minimum size=3.4pt,
                  fill=#1, draw=#1},
    sdot/.default=cgblue,
    sdotb/.style={circle, inner sep=0pt, minimum size=3.4pt,
                  fill=white, draw=#1, line width=0.65pt},
    sdotb/.default=cgblue,
    blochaxis/.style={-{Stealth[length=5pt,width=3.5pt]}, birrengray,
                      line width=0.7pt},
    blochback/.style={birrengray, line width=0.7pt, densely dashed},
  ]

% ─────────────────────────────────────────────────────────────
%  (a)  The Bloch sphere of one spin 1/2
% ─────────────────────────────────────────────────────────────
\node[font=\sf\small\bfseries, text=cgblue] at (0,3.00)
  {(a) \tpcifenglish{one spin: the Bloch sphere}{ein Spin: die Blochkugel}};

% Every named point is placed by the parametrisation, nothing by eye.
\coordinate (O) at (0,0);
\tpcbc{Sup}  {\tpcRb}{0}  {0}      % |up>      = eigenstate of S_3, +
\tpcbc{Sdn}  {\tpcRb}{180}{0}      % |down>    = eigenstate of S_3, -
\tpcbc{Splu} {\tpcRb}{90} {0}      % |+>       = eigenstate of S_1, +
\tpcbc{Smin} {\tpcRb}{90} {180}    % |->       = eigenstate of S_1, -
\tpcbc{Spi}  {\tpcRb}{90} {90}     % |+i>      = eigenstate of S_2, +
\tpcbc{Smi}  {\tpcRb}{90} {270}    % |-i>      = eigenstate of S_2, -
\tpcbc{Spsi} {\tpcRb}{55} {-55}    % a general |psi(theta,phi)>
\tpcbc{Sfoo} {\tpcRb*sin(55)}{90}{-55}   % its foot in the equatorial plane
% Axis tips, all three at the same distance on the page (1.95 cm).
\tpcbc{Axa}  {3.880}{90}{0}
\tpcbc{Axb}  {2.131}{90}{90}
\tpcbc{Axc}  {2.044}{0} {0}
% Far ends of the three axes, just past the far surface.
\tpcbc{Bxa}  {1.05*\tpcRb}{90} {180}
\tpcbc{Bxb}  {1.05*\tpcRb}{90} {270}
\tpcbc{Bxc}  {1.05*\tpcRb}{180}{0}

% Far halves of the axes: dashed, they run behind the ball
\draw[blochback] (O) -- (Bxa);
\draw[blochback] (O) -- (Bxb);
\draw[blochback] (O) -- (Bxc);

% Equator: far (upper) half dashed, near (lower) half solid.  The axis
% ratio 0.3 = sin(eps) is forced by the projection, not chosen.
\draw[birrengray, line width=0.55pt, densely dashed]
  (-\tpcRb,0) arc (180:0:{\tpcRb} and {0.3*\tpcRb});
\draw[birrengray, line width=0.55pt]
  (-\tpcRb,0) arc (180:360:{\tpcRb} and {0.3*\tpcRb});

% Outline
\draw[birrengray, line width=0.9pt] (0,0) circle (\tpcRb);

% Near halves of the axes
\draw[blochaxis] (O) -- (Axa);
\draw[blochaxis] (O) -- (Axb);
\draw[blochaxis] (O) -- (Axc);
\node[lbl, anchor=north east, xshift=-1pt] at (Axa) {$\widehat{S}_1$};
\node[lbl, anchor=west, xshift=2pt]        at (Axb) {$\widehat{S}_2$};
\node[lbl, anchor=south, yshift=1pt]       at (Axc) {$\widehat{S}_3$};

% The general state: its arrow, its foot in the equatorial plane, and the
% two angles of the parametrisation.
\draw[tpc thin, birrengray, densely dotted] (O)    -- (Sfoo);
\draw[tpc thin, birrengray, densely dotted] (Sfoo) -- (Spsi);
\draw[-{Stealth[length=6pt,width=4pt]}, birrenorange, line width=1.1pt]
  (O) -- (Spsi);

\draw[birrengray, line width=0.55pt]
  plot[domain=0:55, samples=30, variable=\s]
    ({\tpcbx{0.62}{\s}{-55}},{\tpcby{0.62}{\s}{-55}});
\node[lbl, font=\scriptsize] at (-0.37,0.66) {$\theta$};

\draw[birrengray, line width=0.55pt]
  plot[domain=0:-55, samples=30, variable=\p]
    ({\tpcbx{0.55}{90}{\p}},{\tpcby{0.55}{90}{\p}});
\node[lbl, font=\scriptsize] at (-0.70,-0.25) {$\varphi$};

% The six eigenstates of S_1, S_2, S_3.  Colour = observable;
% filled = near hemisphere, open = far hemisphere.
\node[sdot=spacecadet]  at (Sup)  {};
\node[sdotb=spacecadet] at (Sdn)  {};
\node[sdot=cgblue]      at (Splu) {};
\node[sdotb=cgblue]     at (Smin) {};
\node[sdot=munsell]     at (Spi)  {};
\node[sdotb=munsell]    at (Smi)  {};

\node[lbl, anchor=west, xshift=3pt]       at (Sup)  {$\ket{\uparrow}$};
\node[lbl, anchor=west, xshift=3pt]       at (Sdn)  {$\ket{\downarrow}$};
\node[lbl, anchor=east, xshift=-4pt]      at (Splu) {$\ket{+}$};
\node[lbl, anchor=south west, xshift=2pt] at (Smin) {$\ket{-}$};
\node[lbl, font=\scriptsize, text=birrengray, anchor=north, yshift=-4pt]
                                          at (Spi)  {$\ket{{+}i}$};
\node[lbl, font=\scriptsize, text=birrengray, anchor=south, yshift=4pt]
                                          at (Smi)  {$\ket{{-}i}$};
\node[lbl, text=birrenorange, anchor=south east, yshift=1pt]
                                          at (Spsi) {$\ket{\psi}$};

% The formula the whole panel is drawn from, and how to read it.
\node[lbl, font=\scriptsize] at (0,-2.38)
  {$\ket{\psi}=\cos\tfrac{\theta}{2}\,\ket{\uparrow}
    +e^{i\varphi}\sin\tfrac{\theta}{2}\,\ket{\downarrow}$};
\node[lbl, font=\scriptsize, text=birrengray] at (0,-2.82)
  {\tpcifenglish{half angle: orthogonal $\Rightarrow$ antipodal}%
                {halber Winkel: orthogonal $\Rightarrow$ antipodal}};

% ─────────────────────────────────────────────────────────────
%  (b)  The same arrow, for one spin of an entangled pair
% ─────────────────────────────────────────────────────────────
\begin{axis}[
    tpcplot,
    at={(4.30cm,-1.70cm)}, anchor=south west,
    width=7.0cm, height=4.9cm,
    clip=false,
    xmin=0, xmax=90, ymin=0, ymax=1.42,
    xtick={0,22.5,45,67.5,90},
    xticklabels={$0$,$\pi/8$,$\pi/4$,$3\pi/8$,$\pi/2$},
    ytick={0,0.5,1},
    xlabel={$\alpha$},
    title={(b) \tpcifenglish{a pair: arrow vs.\ entanglement}%
                            {ein Paar: Pfeil gegen Verschränkung}},
    legend style={font=\scriptsize, draw=light-gray,
                  at={(0.5,-0.26)}, anchor=north, legend columns=2},
]

  % The EPR point alpha = pi/4: a neutral reference rule
  \addplot[levellinegrey] coordinates {(45,0) (45,1)};

  % Length of spin 1's Bloch arrow
  \addplot[cgblue, line width=1.0pt, domain=0:90, samples=181]
    {abs(cos(2*x))};
  \addlegendentry{$\tfrac{2}{\hbar}\bigl|\langle\widehat{\mathbf{S}}\otimes\mathbb{I}\rangle\bigr|=|\cos 2\alpha|$}

  % Failure of the product condition c_11 c_22 = c_12 c_21
  \addplot[birrenorange, line width=1.0pt, densely dashed,
           domain=0:90, samples=181] {abs(sin(2*x))};
  \addlegendentry{$2|c_{11}c_{22}-c_{12}c_{21}|=|\sin 2\alpha|$}

  \addplot[only marks, mark=*, mark size=2pt, cgblue, draw=none]
    coordinates {(45,0)};
  \addplot[only marks, mark=*, mark size=2pt, birrenorange, draw=none]
    coordinates {(45,1)};

  % The family the panel is drawn from
  \node[lbl, font=\scriptsize] at (axis cs:45,1.26)
    {$\ket{\psi_\alpha}=\cos\alpha\,\ket{\uparrow\uparrow}
      +\sin\alpha\,\ket{\downarrow\downarrow}$};

  % What the three marked values of alpha are
  \node[lbl, font=\scriptsize, text=birrengray, anchor=north]
    at (axis cs:0,-0.15)  {$\ket{\uparrow\uparrow}$};
  \node[lbl, font=\scriptsize, text=birrengray, anchor=north]
    at (axis cs:45,-0.15) {$\ket{\text{EPR}}$};
  \node[lbl, font=\scriptsize, text=birrengray, anchor=north]
    at (axis cs:90,-0.15) {$\ket{\downarrow\downarrow}$};
\end{axis}

\end{tikzpicture}
\end{center}


The arrow on the left is the expectation value $\tfrac{2}{\hbar}\langle\widehat{\mathbf{S}}\rangle$ of a single spin. On the right one sees it shrink as soon as the spin is entangled: for the EPR state it has retreated to the centre of the sphere -- there, where the dashed curve has its maximum.

Such states are called \emph{entangled}, and today they are produced routinely in the laboratory.
